\section{Q\&A y recomendaciones prácticas}

\subsection{Puntos principales de la sesión de preguntas}

En la sesión de Q\&A, Chelsea Finn enfatizó los tres pasos «observe → instrument → govern», en particular la necesidad de identificar el prompt, los datos o el attention artifact concretos después de que ocurre un emergent drift.

\begin{table}[H]
\centering
\begin{tabular}{p{4cm}p{5cm}}
\toprule
Tema de la pregunta & Respuesta central \\
\midrule
Emergent stability & Solo después de que varias métricas coinciden puede declararse madura la emergence \\
Prompt drift & insist on prompt log + rollback plan \\
Tooling stack & Las trazas de attention/plan/metric deben ser realtime + archival \\
\bottomrule
\end{tabular}
\caption{Lo esencial de la sesión de Q\&A.}
\end{table}

\begin{knowledgebox}{Puntos para recordar del Q\&A}
Escribe las keywords del Q\&A en `qa.md' y contrástalas en cada incident review, para asegurarte de no omitir ninguna advertencia de alto riesgo.
\end{knowledgebox}

\subsection{Recomendaciones prácticas}

\begin{figure}[H]
\centering
\includegraphics[width=0.85\textwidth]{slides-images/slide-41.jpg}
\caption{Slide de Q\&A y recomendaciones prácticas (slide-41).}
\end{figure}

\begin{importantbox}{Recomendaciones para la puesta en práctica}
\begin{itemize}
    \item Registrar la prompt version + sampling seed;
    \item cada emergent drill debe producir lessons learned + summary note;
    \item en las reuniones del governance panel, hacer review del incident grade-book + prompt log.
\end{itemize}
\end{importantbox}

\subsection{Resumen de la sección}

El Q\&A convierte los abstract insights de la clase en una checklist; mientras los artifacts que surgen de las preguntas y respuestas se sigan poniendo por escrito, la gobernanza siempre tendrá un fundamento verificable.
